\scanpage{111}
\mspage{56}
% archive sheet number: 110
ce qui suppose donc qu'on a choisi des foncteurs image inverse $u_X^*:X_t\longrightarrow X_s$, pour la catégorie $\mathrm{Cat}$-fibrée $X/S$. (Mais \deleted{\unsure{tout}} à \unsure{isom|iso} unique près, \unsure{$\sigma$|ça|$r$|$\gamma$} ne dépend pas de ces choix.) Je me dispense de la vérification fastidieuse que cette loi est bien ``fonctorielle \inserted{en $s$}'', i.e. définit bien une flèche cartésienne

\[
 \sigma_X:\Phi(X)\longrightarrow X,
\]

et de vérifier que celui-ci est fonctoriel en $X$.

Peut-être est-il plus élégant de montrer que $\deleted{\unsure{p}}\pi\inserted{:\underline{\mathrm{Fibsc}}\,S\longrightarrow\mathrm{Cat}_{/S}}$, $\deleted{\unsure{\varphi}}\Phi\inserted{:\mathrm{Cat}_{/S}\longrightarrow\underline{\mathrm{Fibsc}}_{/S}}$ sont adjoints l'un de l'autre :

\[
 \boxed{\mathrm{Hom}\,\mathrm{cart}_S\bigl(\deleted{\varphi}\Phi
 \underset{\substack{\begin{tikzpicture}[baseline=-.3ex]\node[rotate=-90,inner sep=0pt] {$\scriptstyle\in$};\end{tikzpicture}\\\deleted{\unsure{\mathrm{Fib}}}\mathrm{Cat}_{/S}}}{X},
 \underset{\substack{\begin{tikzpicture}[baseline=-.3ex]\node[rotate=-90,inner sep=0pt] {$\scriptstyle\in$};\end{tikzpicture}\\\underline{\mathrm{Fibsc}}}}{Y}\bigr)
 \simeq\mathrm{Hom}_S(X,\deleted{p}\pi Y)},
\]

d'où résulteraient des flèches d'adjonction

\[
 \deleted{\varphi}\Phi\deleted{p}\pi\longrightarrow\mathrm{id}_{\underline{\mathrm{Fibsc}}},
 \qquad \deleted{p}\pi\deleted{\varphi}\Phi\longleftarrow\mathrm{id}_{\underline{\mathrm{Fib}}},
\]

et on aura gagné. Mais pour prouver cette formule d'adjonction, j'ai l'impression qu'il faut \deleted{\unsure{p}} utiliser $\varphi p\mathbin{\unsure{-|\to}}\mathrm{id}$ \ldots\ Il y a des vérifications fastidieuses \ldots

Finalement, je vois maintenant qu'on a \deleted{\unsure{quatre}} \inserted{cinq} cas à traiter, et trois si on joue sur les symétries :

